\section{A seven-generated example}
\label{sec:nonsymmetric-example}

We now consider the non-symmetric numerical semigroup
$$
 H=\langle4,9,15\rangle.
$$
In the notation used for the case $v<2u$ in
Subsection~\ref{subsec:nonsymmetric-multiplicity-four}, this example
corresponds to $(b,c,e)=(6,4,0)$.
Relabel $x_1,x_2,x_3$ as $x,y,z$, respectively, so that
$\deg x=4$, $\deg y=9$, and $\deg z=15$.
Its Ap\'ery set with respect to $4$ is $\{0,9,15,18\}$, and the defining
ideal of $k[H]$ is
$$
 \fkp_H=(x^6-yz,\ y^3-x^3z,\ z^2-x^3y^2).
$$
The degrees of these three relations are $24$, $27$, and $30$, respectively.
Thus the least relation degree is $\ell=24$.  The least elements of $H$ that
are at least $25$ in the four residue classes modulo $4$ are represented by
$$
 x^7,\qquad x^4y,\qquad x^3z,\qquad x^2y^2.
$$
It follows that
$$
 I=I_{25}=(g_1=yz-x^6,\ g_2=x^7,\ g_3=x^4y,\ g_4=x^2y^2,
             \ g_5=x^3z,\ g_6=y^3,\ g_7=z^2).
$$
Indeed, the last two generators are obtained from the relations
$y^3-x^3z$ and $z^2-x^3y^2$.  By
Proposition~\ref{prop:inverse-image-ideals}(1), the ideal $I$ is integrally
closed, and Theorem~\ref{thm:classification}(3) shows that these seven
generators are minimal.

\subsection{Cohen--Macaulayness}

\begin{proposition}
\label{prop:example-local-reduction}
Let
$$
 Q=(g_1,g_4+g_5,g_3+g_6+g_7).
$$
Then
$$
 I^2=QI+\m I^2.
$$
Consequently, the reduction number of $IS_\m$ is exactly one, and
$\calR(I)$ is Cohen--Macaulay.
\end{proposition}

\begin{proof}
Since $I=Q+(g_2,g_3,g_5,g_7)$, it is enough to treat the ten products
$g_ig_j$ with $i,j\in\{2,3,5,7\}$.  We first note that
$g_2g_6=xg_3g_4$, $x^8yz=xg_3g_5$, and
$g_2g_4=xg_3^2$ all belong to $\m I^2$.  Moreover,
$$
 g_5g_6=x^3y^2g_1+g_2g_4\in QI+\m I^2.
$$
The remaining products are treated in the following order.
\begin{itemize}
 \item $g_2g_7=xg_5^2\in\m I^2$.
 \item
 $g_2g_3=(g_3+g_6+g_7)g_2-g_2g_6-g_2g_7\in QI+\m I^2$.
 \item
 $g_5^2=(g_4+g_5)g_5-x^5yg_1-g_2g_3\in QI+\m I^2$.
 \item
 $g_5g_7=(g_4+g_5)g_7-x^2yzg_1-x^8yz\in QI+\m I^2$.
 \item
 $g_3g_5=(g_3+g_6+g_7)g_5-g_5g_6-g_5g_7
 \in QI+\m I^2$.
 \item
 $g_2g_5=(g_4+g_5)g_2-g_2g_4\in QI+\m I^2$.
 \item
 $g_3g_7=x^4zg_1+g_2g_5\in QI+\m I^2$.
 \item
 Since $g_4g_5=x^5yg_1+g_2g_3$, we have
 $g_4^2=(g_4+g_5)g_4-g_4g_5\in QI+\m I^2$.  Hence
 $g_3^2=(g_3+g_6+g_7)g_3-g_4^2-g_3g_7\in QI+\m I^2$.
 \item
 We have $g_6g_7=y^2zg_1+xg_4g_5\in QI+\m I^2$.  Therefore
 $g_7^2=(g_3+g_6+g_7)g_7-g_3g_7-g_6g_7\in QI+\m I^2$.
 \item $g_2^2=xg_3g_5-xg_1g_2\in QI+\m I^2$.
\end{itemize}
Thus $I^2=QI+\m I^2$.

Since $I$ is $\m$-primary, localization at $\m$ preserves the number of
minimal generators.  Hence $IS_\m$ is not a parameter ideal, and its
reduction number cannot be zero.  The equality just proved shows that its
reduction number is one.  Proposition~\ref{prop:local-reduction-criterion}
now shows that $\calR(I)$ is Cohen--Macaulay.
\end{proof}

\subsection{Normality}

\begin{proposition}
\label{prop:example-normality}
The ideal $I$ is normal.
\end{proposition}

\begin{proof}
Set $f=yz$ and $J=I+(f)$.  Then
$$
 J=(yz,x^6,x^4y,x^2y^2,x^3z,y^3,z^2)=I_{24}.
$$
Thus $J$ is an integrally closed monomial ideal generated by at most seven
elements, and hence it is normal by \cite{AtakaMatsuoka2026}.  Moreover,
$$
 xf=xg_1+g_2,\qquad
 yf=yg_1+x^2g_3,\qquad
 zf=zg_1+x^3g_5,
$$
so $\m f\subseteq I$.  We shall also use the identities
$$
 fg_2=g_3g_5,\qquad
 fg_4=g_1g_4+g_3^2,\qquad
 fg_6=g_1g_6+g_3g_4,\qquad
 fg_7=g_1g_7+g_5^2.
$$

We use the setup of Subsection~\ref{subsec:normality} with $\mu=7$
and apply Lemma~\ref{lem:coefficient-separation}.

Suppose that $I^s$ is not integrally closed for some positive integer $s$.
Choose $\xi$, $m$, and $\omega$ as in that subsection, so that
$$
 \xi=\xi_0+f^m\omega,
$$
where
$\xi_0\in I^s+fI^{s-1}+\cdots+f^{m-1}I^{s-m+1}$ and
$\omega\in I^{s-m}$ is homogeneous of degree $\delta-24m$.
The four identities above allow us to remove every term involving
$g_2$, $g_4$, $g_6$, or $g_7$.  We may therefore write
$$
 \omega=\sum_{\lambda\in\Lambda}c_\lambda g^\lambda,
 \qquad c_\lambda\in k,
$$
where $g^\lambda=g_1^{\lambda_1}g_2^{\lambda_2}\cdots g_7^{\lambda_7}$ and
$$
\begin{aligned}
 \Lambda=\{\lambda=(\lambda_1,\lambda_2,\ldots,\lambda_7)\in\mathbb N^7
 \mid {}&\lambda_1+\lambda_2+\cdots+\lambda_7=s-m,\\
 &\lambda_1\deg g_1+\lambda_2\deg g_2+\cdots+\lambda_7\deg g_7
       =\delta-24m,\\
 &\lambda_2=\lambda_4=\lambda_6=\lambda_7=0\}.
\end{aligned}
$$

For $\varepsilon_1,\varepsilon_2\in k^\times$, define a homomorphism
$\psi_{\varepsilon_1,\varepsilon_2}\colon S\to k[[t]]$ by
$$
 x\longmapsto t^2,\qquad
 y\longmapsto\varepsilon_2t^5,\qquad
 z\longmapsto\varepsilon_2^{-1}t^7(1+\varepsilon_1t).
$$
Put $\nu=13$.  The elements occurring in $\omega$ have the following
images:
\[
\begin{array}{c|c|c}
 \text{Element}&\psi_{\varepsilon_1,\varepsilon_2}&
 \ord_t\psi_{\varepsilon_1,\varepsilon_2}\\ \hline
 f&t^{12}(1+\varepsilon_1t)&\nu-1\\
 g_1&\varepsilon_1t^{13}&\nu\\
 g_3&\varepsilon_2t^{13}&\nu\\
 g_5&\varepsilon_2^{-1}t^{13}(1+\varepsilon_1t)&\nu
\end{array}
\]
Moreover,
$$
 \ord_t\psi_{\varepsilon_1,\varepsilon_2}(g_2)
 =\ord_t\psi_{\varepsilon_1,\varepsilon_2}(g_4)
 =\ord_t\psi_{\varepsilon_1,\varepsilon_2}(g_7)=14,
 \qquad \ord_t\psi_{\varepsilon_1,\varepsilon_2}(g_6)=15.
$$
Hence
$\psi_{\varepsilon_1,\varepsilon_2}(I)k[[t]]\subseteq(t^\nu)$.

Set
\[
 G=\{i\mid\ord_t\psi_{\varepsilon_1,\varepsilon_2}(g_i)=\nu\}=\{1,3,5\}.
\]
Every $\lambda\in\Lambda$ has support in $G$. The parameter exponent
map in Lemma~\ref{lem:coefficient-separation} is
$\lambda\mapsto(\lambda_1,\lambda_3-\lambda_5)$.

It remains to check that distinct elements of $\Lambda$ give distinct
parameter monomials.  The exponent of $\varepsilon_1$ determines $\lambda_1$,
and the exponent of $\varepsilon_2$ determines $\lambda_3-\lambda_5$.  On the other hand,
the number of factors gives
$$
 \lambda_3+\lambda_5=s-m-\lambda_1.
$$
These two integer equations determine $\lambda_3$ and $\lambda_5$ uniquely
whenever a solution exists.  Thus every $c_\lambda$ is zero.  This gives
$\omega=0$, contrary to the minimality of $m$.  It follows that $I^s$ is
integrally closed for every $s\geq1$, and hence $I$ is normal.
\end{proof}

Combining Propositions~\ref{prop:example-local-reduction}
and~\ref{prop:example-normality}, we obtain the following.

\begin{theorem}
\label{thm:example-rees}
For the ideal $I=I_{25}$ arising from $H=\langle4,9,15\rangle$, the Rees
algebra $\calR(I)$ is a Cohen--Macaulay normal domain.
\end{theorem}
